\chapter{Station Signaux}

\section{Mandat de la station}
La station Signaux transforme des donnees de marche en informations
predictives evaluees hors echantillon. Son objectif n'est pas de produire une
strategie executable, mais de produire des \textit{findings} robustes
\cite{deprado2018}.

\section{Pipeline A-G}
Le pipeline de recherche est structure en sept couches.
La Figure~\ref{fig:signals-ag} en donne la vue synthetique.

\begin{figure}[H]
\centering
\fbox{\parbox{0.92\textwidth}{
\centering
A Candidats \(\rightarrow\) B Eval OOS \(\rightarrow\) C Robustesse \(\rightarrow\)
D Survivants \(\rightarrow\) E Redondance \(\rightarrow\) F Signal courant
\(\rightarrow\) G Ensemble
}}
\caption{Pipeline A-G de la station Signaux}
\label{fig:signals-ag}
\end{figure}

La Figure~\ref{fig:signals-ag} sert de reference aux sous-sections suivantes.

\section{Familles d'indicateurs et interpretation metier}
La version actuelle utilise quatre familles : SMA (tendance), RSI
(oscillation), MACD (momentum) et OBV (volume). Cette couverture est un compromis
entre diversite informationnelle et maitrise du risque de tests multiples.

\section{Evaluation OOS et score de robustesse}
Chaque candidat est evalue sur des fenetres non vues pendant son ajustement.
Un score de robustesse operationnel est construit a partir de plusieurs
composantes :
\begin{equation}
R_i = w_1 \cdot \text{Sharpe}_{i}^{\text{OOS}} + w_2 \cdot \text{HitRate}_{i}
- w_3 \cdot \sigma_i - w_4 \cdot \text{DD}_{i}
\label{eq:robustness}
\end{equation}
avec $\sum_k w_k = 1$.
L'equation~\ref{eq:robustness} formalise la logique de fiabilite utilisee dans
la phase de selection.

\section{Reduction de redondance}
Apres filtrage par robustesse, la station elimine les variantes trop correlees
afin d'eviter un consensus artificiellement gonfle par des signaux quasi
identiques.

\section{Construction de l'ensemble}
Le signal final par famille est une moyenne ponderee par fiabilite :
\begin{equation}
\text{Score}_{\text{ensemble}} = \frac{\sum_{j=1}^{m} w_j s_j}{\sum_{j=1}^{m} w_j}
\label{eq:ensemble}
\end{equation}
ou $s_j$ est le signal de la variante survivante $j$ et $w_j$ son poids de
robustesse. L'equation~\ref{eq:ensemble} reduit la variance par rapport a une
selection ``winner-takes-all'' \cite{deprado2018}.

\section{Sorties de la station Signaux}
Les sorties principales sont :
\begin{itemize}
\item les scores de consensus par famille et horizon,
\item les variantes representatives et leurs diagnostics,
\item des series indicateurs continues pour l'exploration metier.
\end{itemize}
